\chapter{Quick start}
\label{ch:quickstart}

These two pages take you from nothing to an analyzed simulation of ubiquitin
(\pdb{1UBQ}), a small protein of 76 residues. Each step names the keys to type
and the chapter that explains it; the tutorial in Chapter~\ref{ch:tutorial}
shows every screen. Type each key, then press \key{Enter}. Pressing \key{Enter}
alone accepts the value shown in brackets, such as \prompt{(y)}.

\section*{1. Install ProPrep}

Download the installer for your computer from the releases page,\newline
\url{https://github.com/mjgplab/proprep/releases/latest}.\newline On a Mac with an
Apple chip, open the Terminal and type (replace \texttt{<version>} with the
number in the file name):

\begin{terminal}
bash ~/Downloads/ProPrep-<version>-MacOSX-arm64.sh -b -p ~/ProPrep
\end{terminal}

Intel Macs, Windows and Linux use a different file; see
Chapter~\ref{ch:installation}.

\section*{2. Start ProPrep and make a project}

\begin{terminal}
~/ProPrep/bin/proprep-web
\end{terminal}

ProPrep opens in your web browser (Chapter~\ref{ch:starting}). At
\prompt{Select menu mode [1/2]}, type \key{2} for the Full Menu, which this
manual uses. Press \key{Enter} at \prompt{Save this as your default
preference?}. At \prompt{Project name (.)}, type \typed{ubq}: ProPrep makes a
folder \file{ubq} and saves everything there.

\section*{3. Load ubiquitin}

Type \key{1a}, then \key{1} (Load protein structure), \key{1} (Protein Data
Bank), \key{1} (Enter PDB ID directly). Type \typed{1UBQ}, then \key{1}
(Asymmetric unit). ProPrep downloads the file and shows its title. Type
\key{6} (Done) and then \key{m} to return to the main menu
(Chapter~\ref{ch:structure}).

\section*{4. Build the simulation system}

The Topology Generator writes the two files a simulation needs: the
\term{topology} (\file{.prmtop}, the list of atoms, bonds and force-field
parameters) and the \term{coordinates} (\file{.rst7}, where every atom is). It
also surrounds the protein with water and salt (Chapter~\ref{ch:topology}).

\begin{steps}
  \item Type \key{5a}, then \key{2}. Press \key{Enter} for the protein force
        field (ff19SB) and for modified amino acids (none).
  \item Type \key{1} (None) for DNA, RNA, carbohydrate, lipid and small
        molecule. Then press \key{Enter} at every question up to
        \prompt{Add another salt?}: OPC water, explicit solvent, truncated
        octahedron, 10\,\AA{} buffer, 150\,mM NaCl.
  \item At \prompt{Edit the tLEaP template?} type \key{n}. At
        \prompt{Do you want to continue without bond definitions?} type
        \key{y}; the default \key{n} throws the input away. Press \key{Enter}
        twice more.
  \item Type \key{4} and press \key{Enter}. After a few seconds ProPrep says
        \prompt{All expected topology files were generated successfully!}
\end{steps}

\section*{5. Run a short simulation}

Type \key{m}, then \key{5b} and \key{1} (Chapter~\ref{ch:runningmd}).
\begin{steps}
  \item Type \key{n} (files), \key{1} (predefined protocol), \key{4} (Workshop),
        \key{u} (use as-is), \key{s} (restraints), \key{k} (queue).
  \item At the hardware step type \key{3} (Set all simulations to same
        engine), \key{1} (sander), then \key{n}.
  \item At Review \& Save type \key{3} (run locally now), then \key{1} and
        \key{y}.
\end{steps}

\begin{watchout}
The engine starts as \texttt{pmemd}, which the free installation does not
include. If you skip step~2, the run stops with \prompt{No such file or
directory: 'pmemd'}.
\end{watchout}

ProPrep says this protocol takes about 75~minutes with \texttt{sander}.
Shortened (Chapter~\ref{ch:runningmd}), the same system ran in about a minute
on our test laptop. Type \key{3} in the Execute Simulations menu to watch.

\section*{6. Look at the results}

When every step shows \prompt{Finished}, type \key{4} in the Execute
Simulations menu (from the main menu: \key{5b}, \key{4}; see
Chapter~\ref{ch:analysis}). Press \key{Enter} to pick the most recent
simulation, type \key{2} (Trajectory analysis) and pick the \file{.nc} file by
its number. At \prompt{Select analyses} type \key{1} (RMSD), then \key{1}
(Backbone atoms), and \key{Enter} for frame~0. ProPrep prints the RMSD
statistics and a text plot. Press \key{Enter} twice, type \key{0} to leave, and
\key{x} then \key{y} to exit ProPrep.
